\chapter{Conformal Charges and Ward Identities}
\label{ch:volume-iii-conformal-charges-ward-identities}

\section{Charges as Surface Operators}

Let \(\epsilon^\mu(x)\partial_\mu\) be a conformal Killing vector on flat
Euclidean space, and let
\[
  J^\mu_\epsilon(x)=T^\mu{}_\nu(x)\epsilon^\nu(x)
\]
be the corresponding conserved current at separated points.  If
\(\Sigma\subset\R^D\) is an oriented smooth codimension-one surface with unit
normal \(n_\mu\), define the charge supported on \(\Sigma\) by
\begin{equation}
  Q_\epsilon(\Sigma)
  =
  \int_\Sigma \dd\sigma\,n_\mu T^\mu{}_\nu(x)\epsilon^\nu(x).
  \label{eq:cft-charge-surface}
\end{equation}
Here \(\dd\sigma\) is the induced volume element.  The conservation equation
\(\partial_\mu J^\mu_\epsilon=0\) implies that \(Q_\epsilon(\Sigma)\) is
unchanged under smooth deformations of \(\Sigma\) through regions containing
no operator insertion and no boundary.

This deformation property is the local content of a conserved charge.  A
charge may be evaluated on a large sphere surrounding several insertions, or
on a union of small spheres surrounding them separately.  The equality of
these two evaluations is the Ward identity.

The Lorentzian operator meaning of the same construction is the following.
Let \(\widehat Q_\epsilon(\Sigma)\) be the flux of the conserved current
through a spacelike Cauchy surface \(\Sigma\), with future-directed unit normal.
If a local operator \(\widehat{\mathcal O}(0)\) is inserted between two nearby
Cauchy surfaces \(\Sigma_-\) and \(\Sigma_+\), then
\begin{equation}
  [\widehat Q_\epsilon,\widehat{\mathcal O}(0)]
  =
  \widehat Q_\epsilon(\Sigma_+)\widehat{\mathcal O}(0)
  -
  \widehat{\mathcal O}(0)\widehat Q_\epsilon(\Sigma_-).
  \label{eq:lorentzian-charge-commutator-by-slices}
\end{equation}
Here conservation identifies the charge on \(\Sigma_\pm\) with the same
abstract symmetry generator when no insertion lies between the two surfaces.
Equation \eqref{eq:lorentzian-charge-commutator-by-slices} is the definition
of the commutator after the charge has been represented on the Hilbert spaces
immediately to the future and past of the insertion.  After Euclidean
continuation, the two surfaces close into a small hypersurface surrounding the
insertion, and the commutator becomes the small-surface action used below.

\begin{figure}[H]
  \centering
  \resizebox{0.88\textwidth}{!}{%
  \begin{tikzpicture}[x=1cm,y=1cm,every node/.style={font=\scriptsize}]
    \tikzset{
      slice/.style={green!45!black,line width=0.95pt},
      insertion/.style={circle,fill=violet!80!black,inner sep=1.8pt},
      arrow/.style={-{Latex[length=2mm]},line width=0.8pt}
    }
    \draw[arrow] (-4.7,-1.8)--(-4.7,2.25) node[above] {time};
    \draw[arrow] (-5.0,-1.45)--(3.9,-1.45) node[right] {space};
    \draw[slice]
      (-4.1,-0.65) .. controls (-2.8,-0.35) and (-1.45,-0.85) ..
      (0.0,-0.50) .. controls (1.45,-0.15) and (2.35,-0.55) .. (3.25,-0.25);
    \draw[slice]
      (-4.1,0.95) .. controls (-2.8,1.25) and (-1.45,0.75) ..
      (0.0,1.10) .. controls (1.45,1.45) and (2.35,1.05) .. (3.25,1.35);
    \node[green!45!black] at (3.55,-0.25) {\(\Sigma_-\)};
    \node[green!45!black] at (3.55,1.35) {\(\Sigma_+\)};
    \node[insertion] at (-0.35,0.25) {};
    \node[violet!80!black,below right] at (-0.25,0.19) {\(\widehat{\mathcal O}(0)\)};
    \draw[blue!65!black,dashed,line width=0.8pt]
      (-1.35,-0.78) .. controls (-1.65,-0.05) and (-1.25,0.78) ..
      (-0.42,1.15) .. controls (0.55,1.58) and (1.35,0.75) ..
      (1.22,-0.05) .. controls (1.10,-0.88) and (-0.35,-1.10) ..
      (-1.35,-0.78);
    \node[blue!65!black,align=center] at (1.95,0.35)
      {closed surface\\after continuation};
  \end{tikzpicture}
  }
  \caption{A Lorentzian commutator with a conserved charge is represented by
  the difference between charge insertions on slices immediately after and
  before the local operator.  In Euclidean signature this difference is encoded
  by a closed surface enclosing the insertion.}
  \label{fig:charge-commutator-slices}
\end{figure}

\begin{figure}[H]
  \centering
  \resizebox{0.95\textwidth}{!}{%
  \begin{tikzpicture}[x=1cm,y=1cm,every node/.style={font=\scriptsize}]
    \tikzset{
      insertion/.style={circle,fill=violet!80!black,inner sep=1.8pt},
      surface/.style={green!45!black,line width=0.85pt},
      arrow/.style={-{Latex[length=2mm]},line width=0.8pt}
    }
    \draw[thick] (-5.4,-2.05) rectangle (5.4,2.15);
    \node[anchor=west] at (-5.25,2.45) {\(\R^D\)};

    \draw[surface]
      (-4.35,-1.25) .. controls (-4.85,0.10) and (-3.85,1.85) .. (-1.95,1.75)
      .. controls (-0.85,2.15) and (0.85,2.05) .. (2.05,1.50)
      .. controls (4.05,1.10) and (4.65,-0.65) .. (3.55,-1.40)
      .. controls (2.45,-2.10) and (0.35,-1.65) .. (-1.20,-1.75)
      .. controls (-2.80,-1.85) and (-4.00,-1.80) .. (-4.35,-1.25);
    \node[green!45!black] at (3.95,1.65) {\(\Sigma\)};

    \foreach \x/\y/\lab/\slabel in {
      -2.65/0.55/{x_1}/{S_1},
       0.05/0.85/{x_2}/{S_2},
       2.45/-0.55/{x_n}/{S_n}
    }{
      \node[insertion] at (\x,\y) {};
      \node[violet!80!black,above right] at (\x,\y) {\(\mathcal O(\lab)\)};
      \draw[blue!65!black,line width=0.75pt] (\x,\y) circle (0.55);
      \node[blue!65!black] at (\x+0.52,\y+0.62) {\(\slabel\)};
    }
    \node at (1.35,0.15) {\(\cdots\)};

    \draw[arrow] (5.85,0.55)--(7.15,0.55);
    \node[align=center] at (6.50,1.05) {deform\\surface};

    \begin{scope}[shift={(9.55,0)}]
      \draw[thick] (-2.05,-2.05) rectangle (2.05,2.15);
      \foreach \x/\y/\lab/\slabel in {
        -1.05/0.55/{x_1}/{S_1},
         0.20/0.85/{x_2}/{S_2},
         1.00/-0.55/{x_n}/{S_n}
      }{
        \node[insertion] at (\x,\y) {};
        \node[violet!80!black,above right] at (\x,\y) {\(\mathcal O(\lab)\)};
        \draw[surface] (\x,\y) circle (0.43);
        \node[green!45!black] at (\x+0.40,\y+0.50) {\(\slabel\)};
      }
      \node at (0.68,0.10) {\(\cdots\)};
    \end{scope}
  \end{tikzpicture}
  }
  \caption{A conserved conformal charge can be moved from a large enclosing
  surface to small enclosing surfaces around the insertions.  The Ward identity
  is this equality inside correlation functions.}
  \label{fig:ward-surface-deformation}
\end{figure}

\section{Operator Variations from Small Spheres}

Let \(S_x\) be a small sphere enclosing a single insertion at \(x\), with
orientation inherited from the boundary of a small ball removed from the
correlation-function domain.  The infinitesimal action of the symmetry on the
operator \(\mathcal O(x)\) is defined by
\begin{equation}
  \delta_\epsilon\mathcal O(x)
  =
  \ii\,Q_\epsilon(S_x)\mathcal O(x)
  \label{eq:variation-by-small-charge}
\end{equation}
in Euclidean radial-ordering conventions.  In Lorentzian operator language,
this convention corresponds to
\[
  \delta_\epsilon\widehat{\mathcal O}
  =
  \ii[\widehat Q_\epsilon,\widehat{\mathcal O}].
\]
Changing the orientation convention changes both sides of
\eqref{eq:variation-by-small-charge} by the same sign; the Ward identity
below is then unchanged.

Suppose the vacuum is invariant under the charge, and suppose the large
surface can be moved to infinity with vanishing contribution.  Then for
separated insertions,
\begin{equation}
  0
  =
  \left\langle
    Q_\epsilon(\Sigma)
    \prod_{i=1}^n\mathcal O_i(x_i)
  \right\rangle
  =
  \sum_{i=1}^n
  \left\langle
    \mathcal O_1(x_1)\cdots
    Q_\epsilon(S_i)\mathcal O_i(x_i)\cdots
    \mathcal O_n(x_n)
  \right\rangle ,
  \label{eq:surface-ward-decomposition}
\end{equation}
where \(\Sigma\) encloses all insertions and \(S_i\) is a small sphere
enclosing \(x_i\).  Using
\eqref{eq:variation-by-small-charge}, this becomes
\begin{equation}
  \sum_{i=1}^n
  \left\langle
    \mathcal O_1(x_1)\cdots
    \delta_\epsilon\mathcal O_i(x_i)\cdots
    \mathcal O_n(x_n)
  \right\rangle
  =
  0.
  \label{eq:global-conformal-ward-identity}
\end{equation}
This is the global Ward identity associated with the conformal Killing vector
\(\epsilon^\mu\).  It is a statement about distributions on the configuration
space of distinct points, completed by contact terms on diagonals when
operator products collide.

\section{Source-Derived Local Ward Identities}

The surface identity is the integral form of the source Ward identities
introduced in Section~\ref{sec:cft-source-functionals-ward-contacts}.  Let
\[
  X=\mathcal O_1(x_1)\cdots\mathcal O_n(x_n),
  \qquad x_i\ne x_j ,
\]
where the \(\mathcal O_i\) are scalar local operators.  In the source chart
used for the stress tensor and the operators, translation invariance gives the
distributional identity
\begin{equation}
  \partial_\mu
  \left\langle
    T^{\mu\nu}(x)X
  \right\rangle
  =
  -\sum_{i=1}^n
  \delta^{(D)}(x-x_i)
  \left\langle
    \mathcal O_1(x_1)\cdots
    \partial_i^\nu\mathcal O_i(x_i)\cdots
    \mathcal O_n(x_n)
  \right\rangle .
  \label{eq:local-translation-ward}
\end{equation}
Here \(\partial_i^\nu\) differentiates the \(i\)th insertion coordinate.  A
spin-carrying insertion adds the local representation term dictated by the
diffeomorphism action on its source.  Equation
\eqref{eq:local-translation-ward} is equivalent to the smeared identity
\eqref{eq:cft-smeared-translation-ward}; the displayed delta functions record
the extension of the separated-point distribution across the diagonals
\(x=x_i\).

For scalar primary insertions of dimensions \(\Delta_i\), the flat-space trace
Ward identity in the same source chart is
\begin{equation}
  \left\langle
    T^\mu{}_\mu(x)X
  \right\rangle
  =
  -\sum_{i=1}^n
  \Delta_i\,\delta^{(D)}(x-x_i)
  \left\langle X\right\rangle
  +
  \mathcal A(x;X).
  \label{eq:local-trace-ward}
\end{equation}
The distribution \(\mathcal A(x;X)\) is the local anomaly or contact term
obtained by differentiating the local source functional
\(\mathcal A_\sigma[X]\) in \eqref{eq:cft-smeared-weyl-ward} with respect to
the Weyl parameter \(\sigma(x)\).  It vanishes on the configuration space of
distinct flat-space points; on collision diagonals or curved backgrounds it is
part of the specified source convention.

Let \(\epsilon^\mu\) be a conformal Killing vector, let
\[
  \sigma_\epsilon=\frac1D\partial_\rho\epsilon^\rho,
  \qquad
  J^\mu_\epsilon=T^\mu{}_\nu\epsilon^\nu ,
\]
and assume the stress tensor is symmetric.  Combining
\eqref{eq:local-translation-ward} and \eqref{eq:local-trace-ward} gives
\begin{equation}
  \partial_\mu
  \left\langle J^\mu_\epsilon(x)X\right\rangle
  =
  \sum_{i=1}^n
  \delta^{(D)}(x-x_i)
  \left\langle
    \mathcal O_1(x_1)\cdots
    \delta_\epsilon\mathcal O_i(x_i)\cdots
    \mathcal O_n(x_n)
  \right\rangle
  +
  \sigma_\epsilon(x)\mathcal A(x;X),
  \label{eq:local-conformal-current-ward}
\end{equation}
where for a scalar primary
\[
  \delta_\epsilon\mathcal O_i(x_i)
  =
  -\epsilon^\nu(x_i)\partial_{i,\nu}\mathcal O_i(x_i)
  -\Delta_i\sigma_\epsilon(x_i)\mathcal O_i(x_i).
\]
For spin \(\rho_i\), \(\delta_\epsilon\mathcal O_i\) includes the local
rotation term displayed later in
\eqref{eq:primary-infinitesimal-transform}.  Integrating
\eqref{eq:local-conformal-current-ward} over a region whose boundary is
\(\Sigma\cup_i(-S_i)\) gives the surface decomposition
\eqref{eq:surface-ward-decomposition}; if the large boundary contribution
vanishes and the anomaly term integrates to zero in the chosen background, the
global Ward identity \eqref{eq:global-conformal-ward-identity} follows.

\section{Conformal Charges on the Cylinder}

Under the Weyl map
\[
  \R^D\setminus\{0\}\longrightarrow \R_\tau\times S^{D-1},
  \qquad
  x^\mu=\ee^\tau n^\mu,
  \qquad
  n^\mu n_\mu=1,
\]
closed spheres surrounding the origin become constant-\(\tau\) slices.  The
dilatation charge becomes the cylinder Hamiltonian, rotations become the
rotations of \(S^{D-1}\), and the translation and special conformal charges
become mutually adjoint raising and lowering operators:
\begin{equation}
  \widehat P_\mu^\dagger=\widehat K_\mu .
  \label{eq:p-k-adjoint-cylinder}
\end{equation}
This adjointness is a consequence of reflection positivity under
\(\tau\mapsto-\tau\), which is inversion in flat space.  It is the positivity
input used below to derive unitarity bounds.

More explicitly, let
\[
  \Pi_{\mu\nu}(n)=\delta_{\mu\nu}-n_\mu n_\nu
\]
be the projector onto the tangent space of \(S^{D-1}\) at \(n\).  Acting on
functions on \(\R_\tau\times S^{D-1}\), the flat-space conformal Killing
vectors become
\begin{align}
  D_{\rm rad}
  &=
  \partial_\tau,
  \\
  P_\mu
  &=
  \ee^{-\tau}\left(n_\mu\partial_\tau
  +\Pi_{\mu\nu}(n)\frac{\partial}{\partial n_\nu}\right),
  \\
  K_\mu
  &=
  \ee^\tau\left(-n_\mu\partial_\tau
  +\Pi_{\mu\nu}(n)\frac{\partial}{\partial n_\nu}\right),
  \label{eq:pk-cylinder-vector-fields}
\end{align}
with rotations acting as vector fields tangent to the sphere.  Thus inversion
\(x^\mu\mapsto x^\mu/x^2\), which is \(\tau\mapsto-\tau\), exchanges
\(P_\mu\) and \(K_\mu\).  The factors \(\ee^{-\tau}\) and \(\ee^\tau\) show
that \(P_\mu\) raises and \(K_\mu\) lowers the cylinder energy by one unit in
radial quantization.

The Hamiltonian identification is also visible directly from the stress
tensor.  On the sphere \(S^{D-1}_r\) of radius \(r\), with \(x^\mu=rn^\mu\),
\begin{equation}
  \widehat D_{\rm rad}
  =
  \int_{S^{D-1}_r}\dd\sigma\,n^\mu \widehat T_{\mu\nu}(x)x^\nu
  =
  r^D\int_{S^{D-1}}\dd\Omega\,
  n^\mu n^\nu \widehat T_{\mu\nu}(rn).
  \label{eq:dilatation-charge-stress-tensor-sphere}
\end{equation}
Under the Weyl map this is the cylinder Hamiltonian
\[
  \widehat H_{\rm cyl}
  =
  \int_{S^{D-1}}\dd\Omega\,\widehat T_{\tau\tau}^{\rm cyl}.
\]
In Lorentzian cylinder time \(t=-\ii\tau\), the corresponding Hermitian
generator differs from the Euclidean radial generator by the standard
Wick-rotation factor.  The monograph uses \(\widehat D_{\rm rad}\) for the
positive Euclidean cylinder Hamiltonian that grades local operators by their
scaling dimensions.
